%% 附录

\begin{thesisappendix}

\section{术语与缩写表}

\begin{table}[!h]
    \caption{论文中常用术语与缩写}
    \label{tab:abbreviations}
    \begin{tabular}{p{2.5cm}p{12cm}}
        \toprule
        缩写 / 术语 & 含义 \\
        \midrule
        B-rep        & Boundary Representation，边界表示 \\
        NURBS        & Non-Uniform Rational B-Spline，非均匀有理 B 样条 \\
        CAD          & Computer-Aided Design，计算机辅助设计 \\
        STEP         & ISO 10303 标准 CAD 数据交换格式 \\
        STL          & Stereolithography，三角网格交换格式 \\
        OCC          & OpenCASCADE Community，开源几何建模内核 \\
        VQ-VAE       & Vector-Quantized Variational AutoEncoder，向量量化变分自编码器 \\
        AR           & Autoregressive，自回归 \\
        RoPE         & Rotary Position Embedding，旋转位置编码 \\
        RMSNorm      & Root Mean Square Normalization，均方根归一化 \\
        SwiGLU       & Swish-Gated Linear Unit，门控线性单元 \\
        GQA          & Grouped Query Attention，分组查询注意力 \\
        MHA          & Multi-Head Attention，多头注意力 \\
        MQA          & Multi-Query Attention，多查询注意力 \\
        DDP          & Distributed Data Parallel，分布式数据并行 \\
        AMP          & Automatic Mixed Precision，自动混合精度 \\
        DFS          & Depth-First Search，深度优先搜索 \\
        BRepCheck    & OpenCASCADE 中的 B-rep 有效性校验器 \\
        \bottomrule
    \end{tabular}
\end{table}

\section{代码与数据资源索引}

\begingroup
\small
\begin{sloppypar}
\begin{itemize}
    \item 代码主目录：\path{04-code/nurbs-code/}（核心文件包括 \path{model.py}、\path{trainer.py}、\path{train_ar.py}、\path{generate.py}、\path{sequence2nurbs.py}、\path{utils.py}、\path{eval_generated_brep.py}、\path{dataset.py} 等）。
    \item 数据预处理脚本：\path{04-code/nurbs-code/2brep_nurbs_wcs.py}、\path{04-code/nurbs-code/2brep_nurbs_quan_wcs.py}；服务器侧镜像位于 \path{experiment-local/2sequence_nurbs.py} 与 \path{experiment-local/2brep.py}。
    \item 训练数据：\path{data/nurbs_sequences_v1_patch.pkl}，DeepCAD 子集，训练 64\,325 / 验证 6\,095，数据词表 4\,053。
    \item 实验配置与日志：\path{03-experiments/configs/}、\path{03-experiments/runs/}、\path{03-experiments/metrics/}、\path{03-experiments/toloacl/}（含 \path{hr2_patch1__Val_CE_Loss.csv} 等关键曲线数据）。
    \item 训练日志：\path{experiment-local/logs/hr2_patch1.log}（含模型初始化打印、每 5 epoch 验证最优 checkpoint 记录）。
\end{itemize}
\end{sloppypar}
\endgroup

\section{关键超参数复现表}

下表汇总了主实验 hr2\_patch1 的关键超参数，与正文第~5 章\autoref{tab:model-config} 一致，便于复现。

\begin{table}[!h]
    \caption{hr2\_patch1 关键超参数（复现用）}
    \label{tab:repro-config}
    \begin{tabular}{p{4.5cm}p{9cm}}
        \toprule
        字段 & 取值 \\
        \midrule
        数据集 & DeepCAD 子集 \\
        训练样本数 & 64\,325 \\
        验证样本数 & 6\,095 \\
        数据词表 & 4\,053（50 面索引 + 4\,000 量化坐标 + 3 个结构 token） \\
        模型 \texttt{vocab\_size} & 4\,054（数据词表 + PAD） \\
        \texttt{local\_window\_size} & 1（标准 next-token） \\
        face\_block / edge\_block & 49 / 14 \\
        \texttt{max\_seq\_len} & 4\,096 \\
        $n_{\text{layers}}$ / $n_{\text{heads}}$ / $n_{\text{kv\_heads}}$ & 8 / 8 / 4 \\
        $d_{\text{model}}$ / dim\_ffn & 256 / 1\,024 \\
        RoPE $\theta$ & 500\,000 \\
        RMSNorm $\epsilon$ & $10^{-5}$ \\
        优化器 & AdamW \\
        学习率 / weight decay / dropout & $5 \times 10^{-4}$ / 0.01 / 0.1 \\
        per-GPU batch / 全局 batch & 32 / 128 \\
        训练 epoch 数 & 200 \\
        GPU & 4 $\times$ NVIDIA H20 \\
        采样配置 & top-k = 1, top-p = 0.9 \\
        \bottomrule
    \end{tabular}
\end{table}

\end{thesisappendix}
